\section{Limitations and open problems}\label{sec:limits}

\paragraph{Cutoff, cardinality and seed length.}
Three sizes of a measure $\mu$ on prime powers must be kept apart: its \emph{cutoff} $X=\max\operatorname{supp}\mu$, its
\emph{cardinality} $\abs{\operatorname{supp}\mu}$, and the \emph{seed length} (or entropy) needed to sample it. Our theorems and the
barrier below concern the cutoff certified by the explicit-formula method. A measure with cutoff $X$ has cardinality at most about
$X/\log X$, but thinning or reweighting can make the cardinality much smaller while keeping the cutoff, and a small set can contain
very large primes.

\paragraph{The $(\log q)^2$ barrier.}
Every prime-supported construction analysed through the explicit formula has guaranteed bias
$\asymp\log q/\sqrt X$ in terms of its cutoff $X$. The reason is that about $(\log q)/2\pi$ zeros per unit height can conspire;
spreading the weight over scales changes this only by constants. A measure on primes with cutoff $(\log q)^{2-\delta}$ and bias
$<1$ for all $q$ and all $\chi\ne\chi_0$ would put a prime $\le(\log q)^{2-\delta}$ outside every proper subgroup of $\Gq$, since its
support cannot lie in the kernel of any nontrivial character. That would improve the GRH bound for the least prime outside a
subgroup, which is open \cite{lls2015,cmqr2024}. Montgomery's $\Omega(\log q\log\log q)$ \cite{montgomery1971} limits what is
possible even in principle. Empirically, typical moduli behave like random sets (Section~\ref{sec:experiments}): a random set of
cardinality $\asymp\log q\log X/\varepsilon^2$ is typically $\varepsilon$-biased.

\paragraph{What is formal and what is not.}
Theorems~\ref{thm:A}, \ref{thm:sharp}, \ref{thm:closed} and Corollaries~\ref{cor:A}, \ref{cor:cert} are formal.
Proposition~\ref{prop:opt} is formal only in its abstract core (Lemma~\ref{lem:ranktwo}). Propositions~\ref{prop:convsq}
and~\ref{prop:barrier}, Corollary~\ref{cor:asymp} and Section~\ref{sec:complexity} are informal. Of the convolution squares, only the
family $f_A$ of Theorem~\ref{thm:convclosed} is formalized. Only one numerical instance of Theorem~\ref{thm:closed} is
machine-checked. The closed form loses a factor of about $4$ in support relative to the quadrature constants
(Table~\ref{tab:params}), mostly through the bound on $\mathrm{SP}(a)$ and the decay constants in $I_a(\beta)$.

\paragraph{Trust.}
Our unconditional theorems inherit the trust base of the Lean GRH proof. That proof has passed extensive
machine checks, but at the time of writing no mathematician has read it through, and its companion paper
does not contain a proof of the main inequality \cite{rubinsteinpage}.

\paragraph{Open problems.}
\begin{enumerate}
\item Formalize Propositions~\ref{prop:opt}, \ref{prop:convsq} and~\ref{prop:barrier}, and certify further instances,
  e.g.\ $q\le2^{1024}$ or $\varepsilon=10^{-2}$.
\item Determine the optimal nonnegative kernel on a window of finite length $L$. At $L=4$ it strictly beats convolution squares,
  also among $C^\infty$ kernels and through positivity with a majorant (Propositions~\ref{prop:separation} and~\ref{prop:majorant};
  formal with the constant $\frac98$ for a piecewise-linear kernel, computer-assisted for the sharper constants and the $C^\infty$
  kernel). Certified lower bounds (for instance from the dual problem (EP2) of \cite{cmqr2024}) and the dependence on $L$
  remain open. An explicit bias theorem from such a kernel would need the positivity lemmas for a general majorant
  (Proposition~\ref{prop:majorant}), but no explicit zero counting.
\item Extend to other groups whose characters are Dirichlet characters, such as quotients of $\Gq$ and abelian
  Galois groups over $\mathbb Q$. For nonabelian groups, e.g.\ $\mathrm{Aff}(\Z/q\Z)$, combine with zig-zag or
  semidirect-product arguments.
\item The class-group and isogeny-graph expanders of \cite{jmv2009} need GRH for Hecke $L$-functions of
  imaginary quadratic fields, which remains open.
\end{enumerate}

\paragraph{Acknowledgements and attribution.}
GRH for Dirichlet $L$-functions in Lean is due to Alpöge, with Claude \cite{alpoge2026,grhlean}. The explicit
formula library is \cite{zeta23}; the positivity and gamma-side lemmas and the test function are from
\cite{grhmiller}; the Perron formula is from \cite{schoenfeldlean}; the Borel--Carath\'eodory log-derivative
lemma is from \cite{rhprimeerror}. All are Apache-2.0.

\paragraph{Use of AI.}
This work was carried out with Claude Opus 5.5 (Anthropic), used through Claude Code. Claude proposed and checked
candidate results, wrote the Lean formalization, the numerical experiments and drafts of the text, and searched the
literature. The author chose the problems, directed the work and reviewed the results, and takes responsibility for
all claims. The results marked formal are checked by the Lean kernel, with the axiom audit of
Section~\ref{sec:formal}. Informal results, the exposition and the comparison with the literature rest on the
author's review. Errors found during the work, including an incorrect optimality claim in an earlier draft, are
recorded in the text and in the repository history.
